Table~\ref{tab:court_citations} surveys documented court and commission usage
of each compactness metric.

\begin{table}[ht]
\centering
\caption{Court and commission usage of compactness metrics in redistricting proceedings.}
\label{tab:court_citations}
\begin{tabular}{p{1.2cm}p{4.5cm}p{6.5cm}}
\toprule
Metric & Cases / Proceedings & Usage \\
\midrule
PP &
  \textit{Common Cause v.\ Rucho}, M.D.N.C. (2018);
  \textit{Gill v.\ Whitford}, W.D.Wis.\ (2016) &
  Expert reports cited mean PP for challenged plans; courts cited PP as
  evidence of intentional gerrymandering \\
\addlinespace
Reock &
  \textit{Harper v.\ Hall}, N.C.S.Ct.\ (2022);
  \textit{Gill v.\ Whitford}, W.D.Wis.\ (2016);
  Florida Congressional (2015--16) &
  Expert witnesses used Reock for coastal-boundary-sensitive comparisons;
  cited alongside PP in joint compactness analyses \\
\addlinespace
CH &
  \textit{Shaw v.\ Reno}, 509 U.S.\ 630 (1993);
  \textit{League of Women Voters v.\ Commonwealth}, Pa.\ (2018) &
  Court's ``bizarre shape'' test operationalised via convex hull overlay;
  expert reports used CH to document corridor districts \\
\addlinespace
S &
  Colorado IRC (2021);
  Iowa LSA (ongoing) &
  Colorado commission adopted Schwartzberg-equivalent as primary compactness
  criterion; Iowa LSA uses S in biennial evaluations \\
\addlinespace
LW &
  Illinois Congressional (7th Circuit filings);
  New York State Courts (2022) &
  Expert reports cited LW $> 3$ as elongation threshold; used alongside
  Reock in mixed elongation-circularity profiles \\
\addlinespace
PWC &
  Academic filings (\textit{Fryer-Holden} cited);
  Colorado IRC analytical documents (2021) &
  Used in academic expert reports; Colorado commission included
  population-weighted analysis in its technical appendix \\
\bottomrule
\end{tabular}
\end{table}

\subsection{Key Finding from the Survey}

No metric appears in all cases, and no court has established any single
metric as the exclusive legal standard for compactness.
PP is the most frequently cited metric, appearing in the largest number of
reported expert submissions.
Reock is useful in coastal states where PP's boundary sensitivity is a
litigation vulnerability, but reports should disclose whether the value is exact
MBC Reock or \textsc{bisect}'s centroid-radius proxy.
CH is most cited in ``bizarre shape'' challenges under \textit{Shaw v.\ Reno}.
S is operationally mandated in Colorado.
LW appears in elongation-specific challenges.
PWC is the least established in litigation but is gaining academic traction.

This multi-metric landscape confirms that practitioners need all six metrics,
not any single one.
